% Options for packages loaded elsewhere
\PassOptionsToPackage{unicode}{hyperref}
\PassOptionsToPackage{hyphens}{url}
\PassOptionsToPackage{dvipsnames,svgnames,x11names}{xcolor}
%
\documentclass[
  11pt,
]{article}
\usepackage{amsmath,amssymb}
\usepackage{iftex}
\ifPDFTeX
  \usepackage[T1]{fontenc}
  \usepackage[utf8]{inputenc}
  \usepackage{textcomp} % provide euro and other symbols
\else % if luatex or xetex
  \usepackage{unicode-math} % this also loads fontspec
  \defaultfontfeatures{Scale=MatchLowercase}
  \defaultfontfeatures[\rmfamily]{Ligatures=TeX,Scale=1}
\fi
\usepackage{lmodern}
\ifPDFTeX\else
  % xetex/luatex font selection
\fi
% Use upquote if available, for straight quotes in verbatim environments
\IfFileExists{upquote.sty}{\usepackage{upquote}}{}
\IfFileExists{microtype.sty}{% use microtype if available
  \usepackage[]{microtype}
  \UseMicrotypeSet[protrusion]{basicmath} % disable protrusion for tt fonts
}{}
\makeatletter
\@ifundefined{KOMAClassName}{% if non-KOMA class
  \IfFileExists{parskip.sty}{%
    \usepackage{parskip}
  }{% else
    \setlength{\parindent}{0pt}
    \setlength{\parskip}{6pt plus 2pt minus 1pt}}
}{% if KOMA class
  \KOMAoptions{parskip=half}}
\makeatother
\usepackage{xcolor}
\usepackage[letterpaper,margin=1in]{geometry}
\setlength{\emergencystretch}{3em} % prevent overfull lines
\providecommand{\tightlist}{%
  \setlength{\itemsep}{0pt}\setlength{\parskip}{0pt}}
\setcounter{secnumdepth}{-\maxdimen} % remove section numbering
\usepackage{amsmath,amssymb,amsthm,mathtools}
\usepackage{lmodern,microtype}
\setlength{\emergencystretch}{2em}
\newtheorem{theorem}{Theorem}[section]
\newtheorem{proposition}[theorem]{Proposition}
\newtheorem{lemma}[theorem]{Lemma}
\newtheorem{corollary}[theorem]{Corollary}
\theoremstyle{remark}\newtheorem{remark}[theorem]{Remark}
\newcommand{\R}{\mathbb R}\newcommand{\N}{\mathbb N}\newcommand{\Z}{\mathbb Z}
\newcommand{\norm}[1]{\lVert#1\rVert}\newcommand{\ip}[2]{\langle#1,#2\rangle}
\newcommand{\av}[1]{\langle#1\rangle}\newcommand{\dimA}{\dim_{\mathrm A}}
\newcommand{\dimB}{\mathop{\overline{\dim}}\nolimits_{\mathrm B}}
\DeclareMathOperator{\diam}{diam}\DeclareMathOperator{\spanop}{span}
\ifLuaTeX
  \usepackage{selnolig}  % disable illegal ligatures
\fi
\usepackage{bookmark}
\IfFileExists{xurl.sty}{\usepackage{xurl}}{} % add URL line breaks if available
\urlstyle{same}
\hypersetup{
  pdftitle={Bounded cluster avoidance for uncountable selector families},
  pdfauthor={mxym},
  colorlinks=true,
  linkcolor={blue},
  filecolor={Maroon},
  citecolor={Blue},
  urlcolor={blue},
  pdfcreator={LaTeX via pandoc}}

\title{Bounded cluster avoidance for uncountable selector families}
\author{mxym}
\date{Research note version 1 --- October 7, 2026}

\ifdefined\pdftrailerid\pdftrailerid{}\fi
\ifdefined\pdfinfoomitdate\pdfinfoomitdate=1\fi
\ifdefined\pdfsuppressptexinfo\pdfsuppressptexinfo=15\fi
\begin{document}
\maketitle

\textbf{mxym --- AI-assisted research note, version 1, 7 October 2026.}

A complete written proof with an independent model-conducted
mathematical audit. No human peer review, proof-assistant formalization,
or novelty claim is implied.

\subsection{Result and relation to the existing
entries}\label{result-and-relation-to-the-existing-entries}

This note proves an uncountable-family extension of the routing argument
in mxym/math entries 004 and 006. Instead of making one candidate at a
scale hit a small open set, it makes an entire bounded finite cluster
hit. Every pointwise selector from those clusters is then excluded by
the same large closed set. The proof does not enumerate selectors or
approximate them by a countable family of leading profiles.

Entry 003 concerns a different invariant: upper box dimension zero and
Assouad dimension two of a Banach-space nonembedding obstruction. Its
covering proof suggests keeping branching cardinality separate from
scale count, but its nonembedding theorem is not an analytic dependency
here. In the example below the \emph{space of profiles} has Assouad and
upper box dimension one, whereas every individual null configuration has
upper box dimension zero. Neither of these statements identifies those
dimensions with the logarithmic upper Banach density used in 004. The
critical-covering refinement of 003 is recorded separately in
\href{../critical-covering-gauge/paper.md}{the critical-covering note}.

\subsubsection{Theorem A}\label{theorem-a}

Let S be a subset of the positive integers with \(d^*(S)\)\textgreater0.
For every j in S let \(F_j\) be a nonempty finite set satisfying

\[
F_j\subset[2^{-j-1},2^{-j}],\qquad |F_j|\le K.
\]

where K is a fixed positive integer. Let (\(\omega_r\)) be a nonempty
countable family of nondecreasing finite null moduli, meaning
\(\omega_r\)(a) tends to zero as a tends to zero.

For every \(0<\epsilon<1\) there is a closed nowhere-dense one-periodic
set E such that \(|E\cap[x,x+1]|>1-\epsilon\) for every real x and the
following holds. For every r, y in R, \(c\ne0\), M finite and
nonnegative, and every collection of errors e\_\{j,a\} satisfying

\[
|e_{j,a}|\le M a\omega_r(a).
\]

on all sufficiently late clusters, there are arbitrarily large j in S
for which

\[
\{y+ca+e_{j,a}:a\in F_j\}\subset\mathbb R\setminus E.\tag{A.1}
\]

Consequently, every pointwise selector \(a_j\) in \(F_j\) and every
perturbed sequence \(v_j=y+ca_j+e_j\) with
\(|e_j|\le Ma_j\omega_r(a_j)\) has infinitely many distinct values
outside E on every tail.

The same E works for any prescribed countable family of such cluster
systems; each system may have its own finite bound K and positive
density.

The closed annulus convention is harmless. It permits the useful example
\(F_j=\{2^{-j-1},(3/2)2^{-j-1}\}\) directly. There is no assertion for
unbounded cluster cardinalities or for a continuous interval of
candidates at each scale. The null moduli remain prescribed and
countable.

\subsection{Conventions}\label{conventions}

Here \(d^*(S)\) is upper Banach density: the limit of \(F_S(L)\)/L
defined below. Measure is Lebesgue measure. A configuration has positive
logarithmic upper Banach density when the set of its occupied dyadic
bins has positive upper Banach density. A log-bi-Lipschitz null profile
is a strictly increasing continuous function phi on a small interval
(0,a0{]}, tending to zero at zero, for which
\(\Psi(z)=-\log_2\phi(2^{-z})\) satisfies
\(\alpha(v-u)\le\Psi(v)-\Psi(u)\le\beta(v-u)\) for all sufficiently
large u\textless v, with fixed \(0<\alpha\le\beta<\infty\). Each
prescribed profile may have its own constants. This is the convention of
entry 006 v2.

\subsection{Density template
placement}\label{density-template-placement}

We recall and prove precisely the input from 004/006. Put
\(F_S(L)=\sup_m|S\cap\{m+1,\ldots,m+L\}|\) and \(\delta\)=\(d^*(S)\).
The limit \(F_S(L)\)/L exists by subadditivity and equals its infimum.
For every \(\xi\)\textgreater0 there is a constant \(C_\xi\) such that
every integer interval I obeys

\[
|S\cap I|\le(\delta+\xi)|I|+C_\xi.\tag{1}
\]

Indeed choose \(L_0\) with \(F_S(L_0)\le(\delta+\xi)L_0\) and partition
I into blocks of length \(L_0\) and a remainder; \(C_\xi=L_0\) suffices.
Every tail of S has the same upper Banach density. For every cutoff H
and every length T, some interval of length T beyond H contains at least
\(\delta\) T points. Choose an interval maximizing the count for that
translated tail, rather than requiring it to maximize the original S.
The supremum is attained since the possible counts are integers in a
finite set.

Fix \(0<\eta<\delta\). Suppose an interval template has total span
\(T_L\le C_0L\) and contains finitely many windows, each of length at
least L. Put the template on a length-T\_L interval J beyond any
requested initial index that contains at least \(\delta\) T\_L points,
as supplied by the translated-tail argument. For every window W, its
complement in J consists of at most two intervals. Applying (1) there
gives

\[
|S\cap W|\ge\delta|W|-\xi C_0L-2C_\xi.
\]

Choose \(\xi C_0<(\delta-\eta)/2\) and then L sufficiently large. Every
window contains at least \(\eta\)\textbar W\textbar{} selected indices.
The threshold for L is independent of the requested initial index. This
independence is needed for buffering.

One parity of S has positive upper Banach density; otherwise
subadditivity would force \(d^*(S)\)=0. Replace S by that parity.
Blocking these clusters is enough for Theorem A, and discarding any
finite initial segment is harmless. Hereafter \(\delta\) and \(\eta\)
refer to this chosen parity.

\subsection{Finite routing with entire
clusters}\label{finite-routing-with-entire-clusters}

Fix 0\textless p\textless1/12 and define

\[
\theta=(p/2)^K.
\]

Choose an integer \(M\ge2\) with \(\theta\eta(M-1)>4\). Choose a depth d
with \((1-2^{1-M})^d<p\). Let Q be the number of edges in a complete
ordered M-ary rooted tree of depth d.~Choose an integer gap
g\textgreater=1 with \(Q2^{3-g}<p\).

Order edges in preorder. Edges leaving a vertex of height h have common
window length \(\ell_h\). Put g unused integer indices between
consecutive windows. For a base length L\textgreater=g set

\[
\begin{gathered}\ell_1=L,\qquad\sigma_1=ML+(M-1)g,\\ \ell_h=g+\sigma_{h-1},\\ \sigma_h=M(\ell_h+g+\sigma_{h-1})+(M-1)g\quad(h\ge2).\end{gathered}
\]

The total span \(T=\sigma_d\) is at most \(C_0\) L for fixed M,d,g. If
an edge window is {[}u\_e,v\_e{]} and v\_e* is the final endpoint of it
and its child subtree, then v\_e*-u\_e+1\textless=2 \(\ell_h\). Apply
the density template lemma to place it beyond any initial index, with at
least \(\eta\) \(\ell_h\) selected indices per window. Require the
initial index to be at least four.

Use one-periodic nested grids with \(N_b=2^{b+3}\) cells and key
\(J_b(z)=\lfloor N_b\{z\}\rfloor\). For every nondefault edge e=(P,i),
i\textless M, assign an independent fair-bit selector table at grid
v\_e. The Mth child is default. Every leaf has an independent
Bernoulli-p terminal table at its incoming edge's grid. All entries in
all tables are independent.

Route z to the first nondefault child whose addressed selector is one,
or to the default when all are zero, and continue to a leaf. Let B be
the set of z whose addressed terminal bit is one. Conditioning on all
selectors shows E density(B)=p.~Every outcome is constant on a finest
periodic grid with \(N=2^{u+T+2}\), where u is the template start.

A center x is stable if for each predecessor endpoint b' and next window
start u'=b'+g+1, the interval (x,x+2\^{}(1-u'){]} contains no point of
the grid N\_b'\^{}(-1) Z. The unstable centers have density at most
\(Q2^{3-g}<p\). For a stable x and any candidate a in \(F_j\) with j in
a window, every x+t a, 1\textless=t\textless=2, agrees with x at all
grids of earlier edges. This follows from t a\textless=2\^{}(1-u\_e),
the stability condition and nesting.

Candidates from \emph{different} selected indices in one window have
distinct keys at that window's grid, and avoid the center key. To check
this, if j=j+2, then

\[
\min F_j-\max F_k\ge2^{-j-1}-2^{-k}\ge2^{-j-2}\ge2^{-v_e-2}.
\]

which is twice the cell width. Also every candidate is at distance at
least 2\^{}(-v\_e-1) from the center. All displacements lie in an
interval of length at most 1/8, so circular wraparound creates no
collision. The separation persists at every finer grid. Candidates
\emph{within} one cluster may share keys, and we allow this explicitly.

Fix a stable x and expose its addressed selector entry in every selector
table, but expose no terminal entries. This determines x's route. Its
probability of having no default is \((1-2^{1-M})^d<p\). On an exposure
atom whose first default is at U of height h, test every selected j in
each of U's first M-1 outgoing windows. For an edge e\_i=(U,i), route
each z=x+t a locally from its child U\_i. If every own-edge selector at
that cluster is one and every resulting local terminal entry is one, the
entire cluster x+t \(F_j\) belongs to B. Earlier routing decisions agree
with x; therefore all its candidates really enter U\_i.

\subsubsection{The new probability
estimate}\label{the-new-probability-estimate}

There are \(m\ge\eta(M-1)\ell_h\) tested clusters. For each cluster, let
\(A_j\) be the event that all its own-edge selector entries are one. It
consults at most K distinct entries, so
\(\mathbb P(A_j\mid\text{center exposure})\ge2^{-K}\). Different
clusters consult disjoint own-edge entries, by the separation just
proved (or use different tables). Thus their \(A_j\) events are
independent under the center exposure.

Now condition on \emph{all} selectors. Let \(d_j\) be the number of
distinct local terminal addresses within cluster j; \(1\le d_j\le K\).
Addresses belonging to different clusters are distinct: different
children have disjoint subtrees; different leaves have different tables;
at the same leaf, different index-clusters retain distinct keys at the
finer grid. Conditional on the selectors, the terminal-success events
for different clusters are therefore independent, and the success
probability of cluster j on \(A_j\) is p\^{}(\(d_j\)). The conditional
probability that all tested clusters fail is

\[
\prod_j[1-\mathbf1_{A_j}p^{d_j}]\le\prod_j[1-\mathbf1_{A_j}p^K].
\]

Averaging over selectors and using the independent \(A_j\) events gives

\[
\mathbb P(\text{all clusters fail}\mid\text{center exposure})\le(1-(p/2)^K)^m\le e^{-\theta m}.\tag{2}
\]

This estimate is an inequality, not the singleton engine's exact
identity. Collisions within a cluster can only improve its success
probability. Shared auxiliary selectors do not invalidate the
calculation: they were conditioned on first, and \(d_j\) was bounded
before averaging.

\subsubsection{Uniformity over every normalized
dilation}\label{uniformity-over-every-normalized-dilation}

A candidate in a tested cluster on edge e\_i consults only grids up to
v\_e\_i*. As t traverses {[}1,2{]}, its entire local outcome changes at
at most \(1+2^{2\ell_h+3}\) deterministic boundary values. There are at
most K(M-1)\(\ell_h\) candidates. Represent every boundary, both
endpoints, and one point of every intervening open interval. At most

\[
D_{K,M}(\ell)=4+2K(M-1)\ell(1+2^{2\ell+3}).
\]

representatives suffice. The whole test vector is constant between the
represented boundaries for every table outcome. Thus the probability
that some t has no entire successful cluster is at most

\[
D_{K,M}(\ell)\exp[-\theta\eta(M-1)\ell].\tag{3}
\]

Since \(\theta\eta(M-1)>4>2\log2\), choose L large enough that (3) is
less than p for every integer \(\ell\)\textgreater=L. The parameters
M,d,g,L,T are all chosen before the actual template location. Combining
the no-default cost with (3) gives probability less than 2p at every
stable center. Boundaries are represented individually; failures
supported only at a boundary are retained.

\subsection{Robust normalized blocker}\label{robust-normalized-blocker}

Given a null modulus \(\omega\), let T be the finite span just
constructed. For a possible start u put

\[
\eta_u=\omega(2^{-u})+2^{-u},\qquad r_u=2^{-u}\eta_u>0.
\]

Choose a lower bound on u so late that

\[
4Nr_u=2^{T+4}\eta_u<p,\qquad N=2^{u+T+2}.
\]

This holds at every later start. The density-template lemma still
supplies a filled template there. Every test candidate
a\textless=2\^{}(-u) has \textbar e\_\{j,a\}\textbar\textless=a
\(\omega\)(a)\textless=\(r_u\).

For each random B let \(B_1=B+(-r_u,r_u)\) and \(B_2=B+(-2r_u,2r_u)\) on
the circle. If B has N grid cells, enlarging them gives

\[
\operatorname{dens}(B_2)\le\operatorname{dens}(B)+4Nr_u.
\]

so \(\mathbb E\operatorname{dens}(B_2)<2p\). Let I=S intersect the union
of all edge windows be the deterministic finite test-index set. Define
the exceptional-center set

\[
\begin{split}R=\{x:\ &\exists t\in[1,2]\ \forall j\in I,\\ &\exists a\in F_j\text{ with }x+ta\notin B_1\}.\end{split}
\]

It is closed and periodic: the failure condition is a finite
intersection of finite unions of closed conditions on the compact
rectangle. Its projection onto the circle is closed. By the finite
estimate and unstable center cost,
\(\mathbb E\operatorname{dens}(R)\le3p\). Choose a table outcome with

\[
\operatorname{dens}(B_2)+\operatorname{dens}(R)<5p.
\]

By outer regularity choose an open periodic V containing R with
\(\operatorname{dens}(V)<\operatorname{dens}(R)+p\). Put
\(H=B_2\cup V\); \(\operatorname{dens}(H)<6p\).

If x is not in R, one tested cluster lies entirely in \(B_1\). Any
collection of errors of size at most \(r_u\) puts the entire cluster in
\(B_2\): the first buffer gives a strict margin and the second absorbs
endpoint equality. If x is in R, it lies in V. As j tends to infinity in
S, every candidate in \(F_j\) and every allowed error tends uniformly to
zero. Therefore all sufficiently late perturbed clusters lie in V. This
repairs every center, not merely almost every center.

We have proved: for every x, every t in {[}1,2{]}, and every allowed
independent error array, at least one entire cluster x+tF\_j+errors lies
in H. The successful index for a fixed x,t can be chosen before the
error array: the whole closed uncertainty interval for every candidate
is contained in H. The statement may be applied to any specified tail
and any dyadic dilation of the cluster system.

\subsection{All coefficients and infinitely many
misses}\label{all-coefficients-and-infinitely-many-misses}

For every r, integer k, integer q\textgreater=1 and tail cutoff h, apply
the normalized blocker to clusters 2\^{}k \(F_j\) for j\textgreater=h,
deleting any finitely many nonpositive annular indices. Their annular
indices are shifted by -k and retain positive upper Banach density. Use
the transformed null modulus

\[
\Omega_{r,k,q}(b)=q2^{-k}\omega_r(2^{-k}b).
\]

Choose positive budgets p\_(r,k,q,h)\textless1/12 with
\(12\sum_{r,k,q,h}p_{r,k,q,h}<\epsilon\). Let U be the union of all
blockers and their reflections, and \(E=\mathbb R\setminus U\). This is
closed and one-periodic, and its complement has density less than
\(\epsilon\).

For c\textgreater0 write \(c=2^kt\) with 1\textless=t\textless=2, and
choose q at least the finite error constant M from Theorem A (not the
local routing arity). At \(b=2^ka\) the error bound M a \(\omega_r\)(a)
is at most b Omega\_(r,k,q)(b). For every sufficiently late h, the
component with that tail supplies an entire cluster as in (A.1). Thus
the good cluster indices are unbounded. For c\textless0 use the
reflected blocker and negate the errors.

Every selected v\_j obeys, on sufficiently late clusters,

\[
(|c|/2)a_j\le |v_j-y|\le(3|c|/2)a_j.
\]

Consequently these outside values tend to y but never equal it there,
which guarantees infinitely many distinct outside values on every tail.
Countably many cluster systems can be added as an extra index to the
summable budget.

Finally E has empty interior. Choose once and for all one selector
\(a_j\) from a bounded tail of the first cluster system. If an interval
lay in E, a sufficiently small translated affine copy of this bounded
selector would fit in that interval. The already proved assertion with
zero errors forbids this. A closed set with empty interior is nowhere
dense. Theorem A is proved.

\subsection{An explicit uncountable log-bi-Lipschitz
family}\label{an-explicit-uncountable-log-bi-lipschitz-family}

Let \(A=\{2^{-n}:n\ge1\}\), and for every binary sequence \(\sigma\)
with \(\sigma_n\in\{1,3/2\}\), prescribe

\[
\phi_\sigma(2^{-n})=\sigma_n2^{-n}.
\]

Define \(\Psi_\sigma\) on {[}1,infinity) by linear interpolation of
\(\Psi_\sigma(n)=n-\log_2\sigma_n\), and set
\(\phi_\sigma(a)=2^{-\Psi_\sigma(-\log_2a)}\) for
0\textless a\textless=1/2, with \(\phi_\sigma(0)=0\). Every segment
slope belongs to

\[
[1-\log_2(3/2),\ 1+\log_2(3/2)].
\]

Both endpoints are positive. Hence every \(\phi_\sigma\) is a strictly
increasing log-bi-Lipschitz null profile with these common constants.
Theorem A applied to \(F_{n-1}=\{2^{-n},(3/2)2^{-n}\}\),
n\textgreater=2, produces one E excluding every image y+c
\(\phi_\sigma\)(A), simultaneously for every \(\sigma\), every
\(c\ne0\), and every y. It also permits the displayed prescribed null
perturbations.

This is an uncountable upgrade: \(\sigma\) may change independently at
every scale. A selector with infinitely many changes has no fixed
constant multiple of the identity as its leading profile. More
generally, a countable family of arbitrary fixed profiles cannot cover
all these selectors by relative o(1) approximation, even allowing
separate scalar coefficients. If two selectors were asymptotic to scalar
multiples of the same profile, their quotient on A would converge. That
quotient belongs to \{2/3,1,3/2\}, so it is eventually constant. Apart
from the two constant-tail classes, this means the two selectors
eventually agree. Each eventual- equality class is countable, while
there are continuum many such classes. This is why countable union
bookkeeping in 006 alone does not supply the new simultaneous
quantifiers.

Equip the binary parameter space with d(\(\sigma\),tau)=2\^{}(-m), where
m is the first differing index. It has Assouad dimension and upper box
dimension exactly one, by counting binary cylinders. Under the uniform
function norm on {[}0,1/2{]}, the map \(\sigma\) -\textgreater{}
\(\phi_\sigma\) is bi-Lipschitz:

\[
(1/2)2^{-m}\le\|\phi_\sigma-\phi_\tau\|_\infty\le3\,2^{-m}.
\]

The lower bound is the difference at a=2\^{}(-m). Before z=m-1 the
profiles agree; afterwards both have size at most 3 2\^{}(-m), which
gives the upper bound. Thus the profile family itself has both
dimensions one. Its individual image sequences have upper box dimension
zero, since they decrease geometrically. The density of their
logarithmic bins is positive. Those are three separate statements about
three separate objects.

These profiles are also uniformly bi-Lipschitz in the ordinary variable
away from zero, and extend continuously to zero; on an interpolation
segment interior their derivatives are the logarithmic slope times
\(\phi_\sigma\)(a)/a. The ratio \(\phi_\sigma\)(a)/a lies in
{[}1,3/2{]}, so the derivative lies between \(1-\log_2(3/2)\) and
\((3/2)(1+\log_2(3/2))\). These common bounds imply ordinary
bi-Lipschitz estimates, which extend to zero by continuity. This does
not assert all such profiles are C1 at zero. For non-eventually-constant
\(\sigma\), \(\phi_\sigma(2^{-n})/2^{-n}\) does not converge, so
differentiability at zero already fails.

The theorem therefore does not contradict the known positive embedding
results for arbitrary bi-Lipschitz maps or the C1 endpoint in entry 006.
It excludes this explicit finite-alphabet family, not all maps with the
same distortion constants, all real power exponents, or flat germs.

\subsection{Exact finite certificates}\label{exact-finite-certificates}

A certificate lists finitely many rational clusters, one nonnegative
rational error radius \(r_a\) per candidate, finitely many rational OPEN
intervals repeated with period one, and a rational density budget. It
asserts

\[
\begin{gathered}\forall x\in[0,1]\ \forall t\in[1,2]\ \exists\text{ a listed cluster }F:\\ [x+ta-r_a,x+ta+r_a]\subset H\quad(\forall a\in F).\end{gathered}\tag{C.1}
\]

Failure means choosing one failed candidate from every cluster, and then
a closed complementary gap {[}l,u{]} reachable by an allowed error. For
each such finite choice the constraints are

\[
0\le x\le1,\quad1\le t\le2,\quad l-r_a\le x+ta\le u+r_a.
\]

Exact half-plane clipping in two rational variables decides feasibility.
Segments and singletons are retained, so endpoint-only failures remain.
If feasible, the checker chooses a rational x,t and independently
verifies one legal escaping error in every cluster. If no choice is
feasible, (C.1) holds. The density is computed by merging genuinely
overlapping open intervals; merely touching intervals retain a singleton
gap.

The implementation check\_cluster\_cover.py uses Fraction arithmetic
only. Seven regression cases pass. The supplied computed certificate is

\[
\begin{gathered}F_n=\{2^{-n},(3/2)2^{-n}\},\quad n=1,2,3,4;\\ r_a=a/100;\\ H=(0,4/5)+\mathbb Z,\qquad\text{density budget }4/5.\end{gathered}
\]

It returns valid after nine feasible-polygon states. The first three
clusters alone fail, and an exact counterexample is
\(x=7/15,\ t=1591/900\), with legal escaping candidates 3/4, 1/4 and
3/16. The four-cluster certificate is nontrivial and tests the stronger
all-candidates assertion. It is not a computed small-density witness for
every \(\epsilon\).

Under the bounded-cardinality and positive-annular-density hypotheses of
Theorem A, for computably listed rational clusters and computable
rational radii with \(r_a\)/a uniformly tending to zero by cluster
index, there is also a terminating exact search for normalized
certificates of any rational positive density budget. The relative radii
admit a finite nondecreasing null-modulus majorant on the annular tail,
obtained by taking the supremum of \(r_a\)/a over all candidates below a
given scale. No computable convergence rate is needed for existence. The
robust normalized proof therefore provides an open H of strictly smaller
density. Entire closed uncertainty intervals are contained in H for a
successful cluster, and this is an open condition on x,t. Compactness
supplies a finite subcover by successful clusters. Shrink to finitely
many rational open intervals inside H still covering those compact
uncertainty intervals, and choose rational endpoints so the strict
density budget survives. Enumerate finite rational interval lists and
cluster prefixes and run the exact verifier. Existence proves the search
terminates. No feasible runtime or complexity bound is claimed.

\subsection{Primary source comparison and
scope}\label{primary-source-comparison-and-scope}

\begin{enumerate}
\def\labelenumi{\arabic{enumi}.}
\tightlist
\item
  OpenAI, The geometric case of the Erdos similarity conjecture, family
  084, pinned commit adc7f1241b42e322a6451854ab7e4b4c146bf78a. Its
  theorem is for each fixed geometric ratio. Sections 04-routing and
  05-scales are the source of center exposure, conditioning on all
  selectors, finite dilation entropy, and exceptional-center repair. The
  bounded-cluster probability estimate (2) is the added step here.
  \href{https://github.com/openai/math/blob/adc7f1241b42e322a6451854ab7e4b4c146bf78a/preprints/The-geometric-case-of-the-Erdos-similarity-conjecture-October-5-2026/build/sections/04-routing.tex}{Pinned
  routing source}
  \href{https://github.com/openai/math/blob/adc7f1241b42e322a6451854ab7e4b4c146bf78a/preprints/The-geometric-case-of-the-Erdos-similarity-conjecture-October-5-2026/build/sections/05-scales.tex}{Pinned
  scales source}
\item
  mxym/math entry 004 v1.1 supplies the
  positive-logarithmic-upper-Banach- density template extension; entry
  006 v1 supplies translation-uniform placement and the two-buffer error
  argument; entry 006 v2 treats a prescribed countable family of
  log-bi-Lipschitz profiles. The exact source snapshots are bundled with
  this note and bound to the immutable mxym/math commit
  1c67b4146f9e54291c8a5ca1f42a028bf307629f; 006 v2 explicitly disclaims
  an arbitrary uncountable profile family.
  \href{https://github.com/mxym/math/blob/1c67b4146f9e54291c8a5ca1f42a028bf307629f/preprints/004-log-density-similarity/v1.1/build/main.tex}{004
  v1.1 source}
  \href{https://github.com/mxym/math/blob/1c67b4146f9e54291c8a5ca1f42a028bf307629f/preprints/006-modulus-nonlinear-similarity/v1/main.tex}{006
  v1 source}
  \href{https://github.com/mxym/math/blob/1c67b4146f9e54291c8a5ca1f42a028bf307629f/preprints/006-modulus-nonlinear-similarity/v2/paper.md}{006
  v2 source}
\item
  Feng, Lai, Xiong, Erdos similarity problem via bi-Lipschitz embedding,
  arXiv:2312.01319v1, Theorem 1.1: a geometrically decaying null
  sequence has a bi-Lipschitz image in every positive-measure set, even
  with derivative one at zero. This is consistent with the restricted
  family here, and warns against claiming all bi-Lipschitz or all C1
  images. \href{https://arxiv.org/html/2312.01319v1}{Primary source}
\item
  Bellissard and Julien, Bi-Lipschitz Embedding of Ultrametric Cantor
  Sets into Euclidean Spaces, arXiv:1202.4330v2, gives geometric context
  for counting weighted binary-tree cylinders and finite Assouad
  dimension. The elementary dimension calculation above is proved here
  directly. \href{https://arxiv.org/abs/1202.4330v2}{Primary source}
\end{enumerate}

This is a written proof with exact finite regression checks and an
independent model-conducted mathematical audit. It is not human peer
review, formal verification, or a literature-priority assertion. The
main argument has no numerical optimization dependency. See
\href{PROOF_AUDIT.md}{the proof audit}, \href{README.md}{verification
instructions}, and \href{../avoidance-covering-provenance/README.md}{the
source provenance}.

\subsection{Extension through the countable log-bi-Lipschitz profiles of
006
v2}\label{extension-through-the-countable-log-bi-lipschitz-profiles-of-006-v2}

The theorem also applies to clusters in a fixed-width logarithmic
annulus. More precisely, replace its annulus condition by

\[
F_j\subset[\lambda2^{-j},\Lambda2^{-j}],\qquad0<\lambda\le\Lambda<\infty.
\]

with \(\lambda\) and \(\Lambda\) fixed for the cluster system. First
divide all candidates by \(\Lambda\). Choose an integer q\textgreater=1
with \(2^{-q}\le\lambda/\Lambda\). One residue class of S modulo q+1 has
positive upper Banach density; keep it. Use grids \(N_b=2^{b+q+3}\),
rather than 2\^{}(b+3). Distinct selected index clusters are separated
by at least \(2^{-j-q-1}\); this is at least four grid cells at an edge
endpoint b\textgreater=j, and each candidate avoids the center key. All
other arguments are unchanged with the following explicit replacements:

\[
\begin{aligned}\text{unstable density: }&Q2^{q+3-g};\\ \text{entropy: }&4+2K(M-1)\ell(1+2^{2\ell+q+3});\\ \text{finest grid: }&N=2^{u+T+q+2};\\ \text{buffer cost: }&4Nr_u=2^{T+q+4}\eta_u.\end{aligned}
\]

Choose g so \(Q2^{q+3-g}<p\). The entropy exponential rate is still 2
log 2, so \(\theta\eta(M-1)>4\) suffices. Positive \(\lambda\) and
finite q keep all choices finite; no dependence on the eventual
translation or tail location is introduced. The preliminary division by
\(\Lambda\) is absorbed into the free coefficient and into the
pulled-back null modulus. This proves the wide- annulus version of
Theorem A with exactly the same conclusion.

\subsubsection{Corollary B}\label{corollary-b}

Let (A\_l) be a nonempty countable family of configurations with
positive logarithmic upper Banach density. Let (\(\phi_r\)) be a
prescribed nonempty countable family of log-bi-Lipschitz null profiles
as in 006 v2. Let (\(D_v\)) be a nonempty countable family of finite
nonempty subsets of (0,infinity), and let (\(\omega_w\)) be a prescribed
nonempty countable family of null moduli. For every
\(\epsilon\)\textgreater0 one closed nowhere-dense one-periodic set E of
measure more than 1-\(\epsilon\) in every unit interval has the
following property.

For every choice of l,r,v,w, y in R, \(c\ne0\) and M finite, every
function f on a sufficiently small tail of A\_l satisfying

\[
\begin{gathered}f(a)=y+cd(a)\phi_r(a)+e(a),\\ d(a)\in D_v,\\ |e(a)|\le M\phi_r(a)\omega_w(a).\end{gathered}\tag{B.1}
\]

has infinitely many distinct values outside E on every sufficiently
small tail. The multiplier d(a) may switch arbitrarily at every input
scale; no regularity or enumeration of these switching functions is
assumed.

Proof. By 006 v2 Lemma 2.1, B=\(\phi_r\)(A\_l) has positive logarithmic
upper Banach density. Choose one b\_j=\(\phi_r\)(\(a_j\)) from every
occupied dyadic bin, then use the cluster \(D_v\) b\_j. Its cardinality
is at most \textbar{}\(D_v\)\textbar, and it lies in
\([\min(D_v)2^{-j-1},\max(D_v)2^{-j}]\), so the wide-annulus theorem
applies. Pull \(\omega_w\) back to the b-coordinate by
\(\Omega(b)=\omega_w(\phi_r^{-1}(b))\), extending constantly beyond a
small interval when necessary. It is a null modulus. Put
\(d_0=\min D_v>0\) and define the candidate-coordinate null modulus
\(\Omega_{\mathrm{new}}(u)=\Omega(u/d_0)\). Since d\textgreater=d0, an
error bounded by \(Mb\Omega(b)\) is at most
\((M/d_0)(db)\Omega_{\mathrm{new}}(db)\), exactly the per-candidate
bound required by the cluster theorem. Prescribe this same error at
every candidate d b in a cluster; whole-cluster success then covers the
actual choice d(\(a_j\)).

The countable indices l,r,v,w are absorbed into the summable density
budget. Unbounded successful cluster indices imply b\_j-\textgreater0,
hence \(a_j\)-\textgreater0 because \(\phi_r\) is increasing and tends
to zero. On sufficiently small inputs, \textbar f(a)-y\textbar{} is
between positive constant multiples of \(\phi_r\)(a), since \(D_v\) is
finite and bounded away from zero and \(\omega_w\)(a)-\textgreater0.
Thus the outside hit values tend to y without equalling y, giving
infinitely many distinct values on every tail. This proves the
corollary.

Taking D=\{1\} recovers the relevant profile statement of 006. Taking
D=\{1,3/2\} permits genuinely arbitrary switching with no fixed leading
profile and gives the stronger uncountable simultaneous conclusion.
Countability still applies to the prescribed profiles, alphabets and
error moduli. No simultaneous claim for all positive real exponents
follows.

Additional finite check: using the same candidate formula and relative
error radius a/100, the first six clusters certify H=(0,3/4)+Z, with
exact density 3/4, in fourteen polygon states. This is saved separately
as verification/toy\_cluster\_certificate\_3\_4.json and
verification/toy\_cluster\_result\_3\_4.json. This smaller finite
example also remains a normalized regression certificate, not a computed
certificate for the arbitrarily-small-density theorem.

\end{document}
